\subsection{Hodge、粗拉普拉斯与规范场边界变分}

外微分的幂零关系与联络的曲率作用共同确定一形式场方程，但二者承担不同的约束。

局部张量表示确定了场的分量，但没有确定哪些分量可以独立传播。对带规范冗余的场，还必须把几何上的微分恒等式与动力学约束同时处理。

标量场的谱只需要 $-\Delta$；带自旋的场却生活在非平凡向量丛中，微分算子必须体现联络及曲率。即使是自旋一，也存在至少两种常用的 拉普拉斯算子。若没有弄清楚两者的差别，就容易错误地把普通函数的球谐谱推广成矢量或张量谱。

在定向黎曼流形上，令 $\Omega^p(M)$ 是光滑 $p$ 形式空间，外微分 $\dd:\Omega^p\to\Omega^{p+1}$ 满足 $\dd^2=0$。取 Hodge 星算子 $*$ 与 $L^2$ 内积，外微分的形式伴随记为 $\delta$。这里 $\Delta_H$ 非负；前面的散度型标量算子 $\Delta_g=\operatorname{div}\operatorname{grad}$ 非正，两者在标量上满足 $\Delta_H=-\Delta_g$。在无边界紧致流形上有
\begin{equation}\label{eq:ghs-hodge-laplacian}
 \Delta_{\rm H}=\dd\delta+\delta\dd.
\end{equation}
$\Delta_{\rm H}$ 在 $p$ 形式上非负，因为
$\langle\Delta_{\rm H}\omega,\omega\rangle
 =\norm{\dd\omega}^2+\norm{\delta\omega}^2$。边界存在时这个公式还会出现边界项，故必须补充绝对或相对边界条件。

\begin{theorem}[Hodge 分解]
\label{thm:ghs-hodge-decomposition}
在紧致、无边界、光滑定向黎曼流形上，$L^2$ 的 $p$ 形式有正交分解
\begin{equation}\label{eq:ghs-hodge-direct-sum}
 L^2\Omega^p=\overline{\operatorname{im}\dd}
 \oplus\overline{\operatorname{im}\delta}\oplus\ker\Delta_{\rm H}.
\end{equation}
每个 de Rham 上同调类有唯一调和代表，故 $\dim\ker\Delta_{\rm H}=b_p(M)$，是拓扑不变量。
\end{theorem}
\begin{proof}
$\Delta_{\rm H}$ 是椭圆非负自伴算子，紧致性使其零本征空间有限维，且 Hilbert 空间按核与正谱正交分解。对正谱部分设 $G=\Delta_{\rm H}^{-1}$，则
\[
 \omega=\mathcal H\omega+\Delta_{\rm H}G(1-\mathcal H)\omega
 =\mathcal H\omega+\dd\delta G\omega+\delta\dd G\omega,
\]
$\mathcal H$ 是调和投影。$\dd^2=\delta^2=0$ 和积分分部说明后两个子空间正交；调和形式与它们也正交。若 $\omega$ 闭合，则 $\dd\omega=0$，上式中的余微分项为零，剩下 $\omega=\mathcal H\omega+\dd\alpha$，说明每个闭形式与唯一的调和形式同调。调和且恰当的形式由范数恒等式只能为零，所以代表唯一。
\end{proof}

这个结论对规范场特别重要：零模个数可以由空间拓扑决定，而不由局域曲率常数决定。例如环面上的一形式具有非平凡调和模式，二球上的一般一形式却没有相同的同调零模结构。

联络导出粗 拉普拉斯算子 $\nabla^*\nabla$。对一形式 $\omega_a$，Weitzenböck 恒等式为
\begin{equation}\label{eq:ghs-weitzenbock-one-form}
 (\Delta_{\rm H}\omega)_a=(\nabla^*\nabla\omega)_a
 +\operatorname{Ric}_a{}^b\omega_b.
\end{equation}
它表明曲率以代数势项的形式改变了自旋一的传播算子。对高阶张量也存在相应曲率作用，但不能只用标量曲率 $\mathcal R$ 乘单位矩阵代替全部 Riemann 张量缩并。

\begin{proof}[证明（一形式曲率项的计算）]
对一形式展开协变导数，并交换指标，可核对曲率项的系数。
取 $\delta\omega=-\nabla^a\omega_a$，则
\begin{align*}
(\dd\delta\omega)_b&=-\nabla_b\nabla^a\omega_a,\\
(\delta\dd\omega)_b&=-\nabla^a(\nabla_a\omega_b-\nabla_b\omega_a)\\
&=-\nabla^a\nabla_a\omega_b+\nabla^a\nabla_b\omega_a.
\end{align*}
相加后得到 $-\nabla^a\nabla_a\omega_b+[\nabla^a,\nabla_b]\omega_a$。按照前文固定的曲率约定，该对易项收缩为 $\operatorname{Ric}_b{}^a\omega_a$，于是得 \cref{eq:ghs-weitzenbock-one-form}。这是带自旋算子与标量 拉普拉斯算子的第一个明确差异。

\end{proof}

考虑 $U(1)$ 规范场 $A\in\Omega^1(M)$，曲率 $F=\dd A$，作用量 $\mathscr I[A]=\frac14\int F_{ab}F^{ab}\dd\mu_g=\frac12\norm{\dd A}^2$。规范变换 $A\mapsto A+\dd\alpha$ 不改变 $F$。在 Lorenz 型规范 $\delta A=0$ 下，形式动力学算子与 $\Delta_{\rm H}$ 密切相关。然而零模可能包括精确形式的纯规范方向，以及非精确调和一形式的全局模态；二者必须区别商掉。

若区域有边界，允许的 $\alpha$ 是在边界消失还是可取非零值，直接决定边界荷与 Hilbert 空间分解。不能仅由局域的自旋一张量表示得出规范场全部物理自由度。高自旋规范场也存在类似的约束、迹条件与背景曲率一致性问题。

对对称二张量 $h_{ab}$，取四指标排列 $R_{abcd}=g(R(e_a,e_b)e_c,e_d)$，并记 $R_a{}^c=\operatorname{Ric}_a{}^c$。可以构造一种二阶椭圆算子，其主符号部分是 $-\nabla^2h_{ab}$，曲率作用为
\begin{equation}\label{eq:ghs-lichnerowicz-two-tensor}
 (\Delta_{\rm L}h)_{ab}=-\nabla^c\nabla_ch_{ab}
 +R_a{}^ch_{cb}+R_b{}^ch_{ac}+2R_{acbd}h^{cd}.
\end{equation}
这里的四指标顺序与前面的 $R(X,Y,Z,W)$ 相同；若先降低 $R^a{}_{bcd}$ 的首指标，所得四指标张量等于本式 $R_{abcd}$ 的负值，曲率项便写成常见的 $-2$ 形式。
\begin{proof}[证明（度量张量的零模核对）]
由度规相容性，$\nabla g=0$。由 Ricci 定义与曲率反对称性，
\[
 g^{cd}R_{acbd}=-\operatorname{Ric}_{ab}.
\]
因此取 $h=g$ 时
\[
 (\Delta_{\rm L}g)_{ab}
 =0+\operatorname{Ric}_{ab}+\operatorname{Ric}_{ab}
 +2(-\operatorname{Ric}_{ab})=0.
\]
常截面曲率 $k$ 的直接代入也给出 $R_{acbd}=k(g_{cb}g_{ad}-g_{ab}g_{cd})$，故其收缩为 $-(d-1)kg_{ab}$，与同一 Ricci 定义相符。
\end{proof}
该算子经常出现在 Einstein 方程的线性化、引力子涨落与相关规范固定问题中。它与单独作用于每个分量的标量 拉普拉斯算子 不是同一个算子。

对于更高阶对称无迹张量，二阶算子的可能曲率修正更多；若还要求规范不变性、幺正性和正确约束传播，允许的背景与耦合受到额外限制。这说明“空间任意”与“自旋任意”是两个独立的推广目标，不能在不增加物理假设的情况下同时任意化。

前面的高自旋微分算子涉及协变导数对易子，但对自旋一场还存在更强的结构：规范变换并不要求曲率消失，它直接来自外微分的恒等式 $\dd^2=0$。这一事实提供了区分“曲率使平直表达式失效”和“曲率破坏规范对称性”的一个可完全算通的例子。

先在无边界紧致定向黎曼流形上讨论欧氏二次型。令实一形式 $A=A_a\dd x^a$ 为规范势，$F=\dd A$ 为二形式场强。由于 Levi--Civita 联络无挠，
\begin{equation}\label{eq:ghs-r9-maxwell-fieldstrength}
 F_{ab}=\nabla_aA_b-\nabla_bA_a
 =\partial_aA_b-\partial_bA_a.
\end{equation}
定义作用量 $S[A]=\tfrac12\int_X F\wedge *F=\tfrac14\int_XF_{ab}F^{ab}\dd\mu_g$。变分时 $\delta F=\dd(\delta A)$，对闭流形积分分部：
\begin{align*}
 \delta S
 &=\int_X\dd(\delta A)\wedge *F
 =\langle F,\dd(\delta A)\rangle_{L^2}\\
 &=\langle\delta F,\delta A\rangle_{L^2}.
\end{align*}
这里最后一个 $\delta$ 在 $\delta A$ 中是变分符号，在 $\delta F$ 中则表示余微分；二者的数学角色不同。Euler--Lagrange 方程为
\begin{equation}\label{eq:ghs-r9-maxwell-coexact-eom}
 \delta\dd A=0.
\end{equation}
对任意光滑标量 $\chi$，取 $A\mapsto A+\dd\chi$，由于 $\dd^2\chi=0$，场强和作用量都不变。这一证明对所有光滑黎曼背景成立，不要求 Ricci 曲率为零。

\begin{proposition}[规范场的一形式 Weitzenbock 恒等式]
\label{prop:ghs-r9-maxwell-weitzenbock}
对一形式 $A$，Hodge 拉普拉斯算子满足
\begin{equation}\label{eq:ghs-r9-maxwell-curvature-term}
 (\dd\delta+\delta\dd)A
 =\nabla^*\nabla A+\operatorname{Ric}(A).
\end{equation}
选取 Lorenz 型规范 $\delta A=0$ 后，自由 Maxwell 方程成为 $\Delta_HA=0$。因此在曲空间直接用粗 拉普拉斯算子 $\nabla^*\nabla$ 代替 $\Delta_H$ 会漏掉 Ricci 项；但该项不是外加的规范对称性破缺项，而是由微分复形决定的几何补偿。
\end{proposition}
\begin{proof}
取正定 Laplace 约定，$\delta A=-\nabla^bA_b$，以及 $F_{ba}=\nabla_bA_a-\nabla_aA_b$。于是
\begin{align*}
 (\delta\dd A)_a
 &=-\nabla^bF_{ba}\\
 &=-\nabla^b\nabla_bA_a+\nabla^b\nabla_aA_b.
\end{align*}
另一方面，$(\dd\delta A)_a=-\nabla_a\nabla^bA_b$，所以
\[
 (\Delta_HA)_a=-\nabla^b\nabla_bA_a
 +[\nabla^b,\nabla_a]A_b.
\]
按本书采用的 Ricci 曲率约定，最后的缩并交换子恰为 $\operatorname{Ric}_a{}^bA_b$。第一项是粗 拉普拉斯算子，这就得到了 \cref{eq:ghs-r9-maxwell-curvature-term}。对 $A=\dd\chi$，还可用 $\Delta_H\dd\chi=\dd\Delta_H\chi$ 检查 Ricci 项准确抵消粗 拉普拉斯算子 与梯度不交换的部分。
\end{proof}

这个推导也解释了矢量球谐中两个容易混用的本征值。单位二球若 $\Delta_HY_{\ell m}=\ell(\ell+1)Y_{\ell m}$，则对 $\ell\ge1$，
\[
 \Delta_H(\dd Y_{\ell m})=\ell(\ell+1)\dd Y_{\ell m},\qquad
 \nabla^*\nabla(\dd Y_{\ell m})=[\ell(\ell+1)-1]\dd Y_{\ell m}.
\]
差值恰是 $S^2$ 的 Ricci 本征值 $1$。在环面上 Ricci 为零，两算子则具有相同本征值。这是上述 Hodge 分解、曲率算子以及表示谱的同一个交叉检验，而不是三个互不相干的公式。

如果 $X$ 有边界，$\delta$ 是 $\dd$ 的形式伴随时会出现边界配对；选择绝对或相对边界条件不仅影响自伴性，也影响允许的规范参数与调和零模。即使无边界，零模 $h$ 满足 $\dd h=\delta h=0$ 时仍可出现在 $A$ 中，且不一定是全局的 $\dd\chi$。例如环面上的 $\dd x,\dd y$ 对应不可用单值周期规范函数消去的平坦连接；在耦合带电场时，还需进一步考察它们的整体 holonomy。

将作用量分部积分时保留边界项，可直接得到相应的条件。
设 $X$ 有光滑边界，$\mathbf n$ 为外法向。对紧支于边界附近但边界取值不一定为零的变分，Stokes 定理给出
\[
 \int_X\langle F,\dd(\delta A)\rangle\dd\mu
 =\int_X\langle\delta F,\delta A\rangle\dd\mu
 +\int_{\partial X}\mathbf n^aF_{ab}\,\delta A^b\dd\mu_{\partial X}.
\]
因此固定切向规范势（Dirichlet 型）时边界变分消失；若允许边界 $A$ 自由变化，则自然条件是相应切向分量的 $\mathbf n^aF_{ab}=0$。一般相对/绝对 Hodge 条件还需同时约束规范固定算子的域，不能只规定 $F$ 的一个分量便宣称椭圆自伴问题已经完全确定。边界上的规范变换还可能成为真实的边缘自由度，因而与无边界的情形存在物理差别。
